\clearpage
\begingroup
\let\clearpage\relax
\let\cleardoublepage\relax
\vspace*{\fill}
\chapter{进程与线程}
\begin{center}
进程是操作系统进行资源分配和调度的基本单位，线程是 CPU 调度的基本单位。本章从进程的状态模型出发，讲清 PCB、上下文切换与线程模型。
\end{center}
\vspace*{\fill}
\endgroup
\clearpage

\section{进程概念}

\noindent 操作系统要在单个 CPU 上同时支撑多个程序的运行，就必须为「正在执行的程序」建立一个统一的抽象与数据结构。这个抽象就是\textbf{进程 (process)}：它既描述了「一段正在运行的代码及其数据」，又是操作系统进行\textbf{资源分配与保护}的独立单元。本节先厘清程序与进程的区别，再说明进程由哪些部分构成，最后区分并发与并行这两个常被混淆的概念。

\subsection{程序与进程}

\noindent \textbf{程序 (program)} 是存放在磁盘上的静态文件，由指令和数据构成，本身不占用 CPU，也不会「运行」。\textbf{进程 (process)} 则是程序在某个数据集合上的一次执行过程，是动态的：它有开始、有状态、有结束，并随执行推进而不断变化。同一个程序可以被多次加载，产生多个相互独立的进程（例如同时打开两个文本编辑器）。

\begin{table}[H]
\centering
\caption{程序与进程的区别}
\begin{tabulary}{\textwidth}{LLL}
\toprule
维度 & 程序 & 进程 \\
\midrule
存在形式 & 磁盘上的静态文件 & 内存中动态执行的实体 \\
组成 & 指令 + 数据 & 代码段 + 数据段 + PCB \\
生命周期 & 长期保存，可反复使用 & 有创建、运行、结束的过程 \\
对应关系 & 一个程序可对应多个进程 & 一个进程通常对应一个（或换入的）程序 \\
资源占用 & 不占 CPU、不占运行时资源 & 占用 CPU、内存、文件等资源 \\
\bottomrule
\end{tabulary}
\end{table}

\noindent 关键结论：程序是\textbf{静态}的，进程是\textbf{动态}的；程序可长期保存，进程有生命周期；同一程序可对应多个进程，而一个进程在运行中还可以通过 \texttt{exec} 换成另一个程序。进程状态用 $\text{state}(p)$ 表示，其生命周期可写作一个状态序列，详见后文。

\subsection{进程的组成}

\noindent 进程在内存中的实体由三部分构成，缺一不可：

\begin{equation}
    \text{进程} = \text{代码段} + \text{数据段} + \text{进程控制块 (PCB)}
\end{equation}

\begin{itemize}
    \item \textbf{代码段 (text)}：可执行的机器指令，通常只读，可被运行同一程序的多进程共享。
    \item \textbf{数据段 (data)}：全局/静态变量、堆与栈等可读写数据，是进程私有内容。
    \item \textbf{PCB}：操作系统为该进程建立的管理数据结构，记录其标识、状态与资源，是进程存在的\textbf{唯一标志}。
\end{itemize}

\noindent 代码段与数据段合称进程的\textbf{地址空间}，PCB 则游离于地址空间之外、由内核统一管理。下图给出进程内存映像与 PCB 的关系。

\begin{figure}[H]
\centering
\begin{tikzpicture}[font=\small, >=stealth, every node/.style={align=center}]
    \draw[rounded corners, blue!35, fill=blue!3] ({0},{0}) rectangle ({3.2},{4.1});
    \node[font=\scriptsize, blue!60!black] at ({1.6},{3.85}) {内存映像};
    \node[draw, fill=green!12, minimum width=2.6cm, minimum height=0.62cm] (stack) at ({1.6},{3.25}) {栈 stack};
    \node[draw, fill=green!12, minimum width=2.6cm, minimum height=0.62cm] (heap) at ({1.6},{2.5}) {堆 heap};
    \node[draw, fill=green!12, minimum width=2.6cm, minimum height=0.62cm] (data) at ({1.6},{1.75}) {数据段 data};
    \node[draw, fill=green!12, minimum width=2.6cm, minimum height=0.62cm] (code) at ({1.6},{1.0}) {代码段 text};

    \node[draw, rounded corners, fill=orange!15, minimum width=3.0cm, minimum height=1.5cm] (pcb) at ({7.0},{2.1}) {PCB\\[2pt]\scriptsize 标识 / 现场 / 调度\\\scriptsize 内存 / 资源 / 记账};

    \draw[<->] ({3.2},{2.1}) -- ({5.5},{2.1}) node[midway, above, font=\scriptsize]{内核管理};
\end{tikzpicture}
\caption{进程的组成：代码段、数据段构成地址空间，PCB 由内核维护}
\end{figure}

\subsection{并发与并行}

\noindent 这是两个极易混淆的概念，务必区分清楚。

\begin{table}[H]
\centering
\caption{并发与并行的区别}
\begin{tabulary}{\textwidth}{LLL}
\toprule
维度 & 并发 (concurrency) & 并行 (parallelism) \\
\midrule
时刻关系 & 宏观同时、微观交替 & 同一时刻真正同时执行 \\
硬件要求 & 单核即可实现 & 必须多核/多处理机 \\
CPU 占用 & 任一时刻只有一个进程在占用 CPU & 多个进程同时占用不同 CPU \\
典型效果 & 提高资源利用率、改善响应 & 提高单位时间完成的计算量 \\
\bottomrule
\end{tabulary}
\end{table}

\noindent 换言之，并发是「进程在时间上交错推进」的能力，并行是「进程真的同时跑」的能力。单核系统只能并发，多核系统才可能既并发又并行。

\section{进程状态模型}

\noindent 进程在生命周期中并非一直在运行：它在等待 CPU、等待 I/O 或等待某个事件。操作系统用一组\textbf{状态}刻画进程此刻的处境，并规定状态之间哪些转换是合法的。状态模型是理解调度与阻塞的关键。

\subsection{五种状态}

\begin{itemize}
    \item \textbf{新建 (new)}：进程正在被创建，PCB 已分配但资源尚未完全就绪，尚未进入就绪队列。
    \item \textbf{就绪 (ready)}：已获得除 CPU 之外的全部资源，只等被调度程序选中即可运行。
    \item \textbf{运行 (running)}：已占有 CPU，正在执行指令（单核上至多一个进程处于此态）。
    \item \textbf{阻塞 (blocked)}：因等待 I/O 或某事件而暂停，即使 CPU 空闲也\textbf{不能}运行。
    \item \textbf{终止 (terminated)}：执行完毕或被撤销，等待系统回收其资源。
\end{itemize}

\subsection{状态转换与触发事件}

\noindent 状态之间的每一次跳转都由一个明确的\textbf{触发事件}引起。合法转换与事件如下表。

\begin{table}[H]
\centering
\caption{进程状态转换及其触发事件}
\begin{tabulary}{\textwidth}{LLL}
\toprule
转换 & 触发事件 & 说明 \\
\midrule
新建 $\to$ 就绪 & 接纳 (admit) & 资源就绪，插入就绪队列 \\
就绪 $\to$ 运行 & 调度 (dispatch) & 被调度程序选中并分派 CPU \\
运行 $\to$ 就绪 & 时间片到 / 被抢占 & 主动或被迫让出 CPU，重新排队 \\
运行 $\to$ 阻塞 & 等待事件（如 I/O 请求） & 进程主动放弃 CPU \\
阻塞 $\to$ 就绪 & 事件完成（如 I/O 完成） & 回到就绪队列，仍不能直接运行 \\
运行 $\to$ 终止 & 退出 (exit) & 正常结束或被强制终止 \\
\bottomrule
\end{tabulary}
\end{table}

\noindent 关键约束：\textbf{阻塞态不能直接回到运行态}，必须先变为就绪态，再由调度程序决定何时运行；就绪到运行由调度触发，运行到就绪由时间片耗尽或被抢占触发。

\begin{figure}[H]
\centering
\begin{tikzpicture}[font=\small, >=stealth, node distance=2.7cm,
    state/.style={draw, rounded corners, align=center, minimum width=1.9cm, minimum height=1.0cm}]
    \node[state, fill=blue!8] (new) {新建\\new};
    \node[state, fill=green!12, right=of new] (ready) {就绪\\ready};
    \node[state, fill=orange!18, right=of ready] (run) {运行\\running};
    \node[state, fill=gray!15, right=of run] (term) {终止\\terminated};
    \node[state, fill=red!12, below=2.5cm of run] (block) {阻塞\\blocked};

    \draw[->] (new) -- node[above, align=center]{接纳\\admit} (ready);
    \draw[->] (ready) -- node[above, align=center]{调度\\dispatch} (run);
    \draw[->] (run) -- node[above, align=center]{退出\\exit} (term);
    \draw[->] (run) to[bend left=45] node[above, align=center]{时间片到\\或抢占} (ready);
    \draw[->] (run) -- node[right, align=center]{等待事件\\I/O 请求} (block);
    \draw[->] (block) -- node[below, align=center]{事件完成\\I/O 完成} (ready);
\end{tikzpicture}
\caption{进程五状态转换图：阻塞态必须先回到就绪态，不能直接进入运行态}
\end{figure}

\subsection{状态转换的方程表示}

\noindent 用一个状态序列可以精确描述进程的一次完整经历。任意合法生命周期都可由如下「骨架」生成：

\begin{equation}
    \text{new} \to \text{ready} \to \text{running}
    \to \{\text{ready},\ \text{blocked},\ \text{terminated}\}
\end{equation}

\begin{equation}
    \text{blocked} \to \text{ready} \to \text{running},
    \qquad
    \text{blocked} \not\to \text{running}
\end{equation}

\noindent 例如某进程 $P$ 的一次执行可记为
$\text{new} \to \text{ready} \to \text{running} \to \text{blocked}
\to \text{ready} \to \text{running} \to \text{ready}
\to \text{running} \to \text{terminated}$。
其中第二段 $\text{blocked} \to \text{ready}$ 由 I/O 完成触发，第三段 $\text{running} \to \text{ready}$ 由时间片耗尽触发。式 (3) 表明：若进程处于阻塞态，它\textbf{不可能}（用 $\not\to$ 表示）直接跳到运行态，忽略这一点是初学时的常见错误。

\section{进程控制块 PCB}

\noindent PCB 是进程存在的\textbf{唯一标志}：创建进程即建立 PCB，撤销进程即回收 PCB。操作系统通过 PCB 感知、控制和调度进程，因此 PCB 中必须完整记录「这个进程是谁、此刻在哪、拥有什么、该怎么调度」。

\begin{figure}[H]
\centering
\begin{tikzpicture}[font=\small, >=stealth]
    \node[draw, rounded corners, fill=blue!10, align=center, minimum width=3.2cm] (title) {进程控制块 (PCB)};
    \node[draw, rounded corners, align=left, inner sep=8pt, fill=gray!4,
          below=0.3cm of title] (pcb) {%
        \begin{tabular}{@{}l@{}}
        进程标识：PID、父进程 PPID、用户 UID\\
        处理机状态：PC、通用寄存器、PSW、栈指针\\
        调度信息：状态、优先级、时间片、队列指针\\
        内存管理：页表基址、段表、界限寄存器\\
        资源占用：打开文件表、设备表\\
        记账信息：CPU 时间、缺页次数、启动时间\\
        同步通信：信号量、消息队列指针\\
        \end{tabular}};
    \draw[decorate, decoration={brace, amplitude=5pt, mirror}]
        ({pcb.south east}) -- ({pcb.north east})
        node[midway, right=7pt, align=left, font=\scriptsize]{进程存在的\\唯一标志};
\end{tikzpicture}
\caption{PCB 的典型字段组织：标识、现场、调度、内存、资源、记账、通信}
\end{figure}

\noindent 各字段的用途逐项说明如下。

\begin{table}[H]
\centering
\caption{PCB 主要字段及其用途}
\begin{tabulary}{\textwidth}{LLL}
\toprule
分类 & 字段 & 用途 \\
\midrule
进程标识 & PID、父进程 PPID、用户 UID & 唯一标识进程及其归属与权限 \\
处理机状态 & 通用寄存器、PC、PSW、栈指针 & 保存/恢复现场，支撑上下文切换 \\
进程调度 & 状态、优先级、时间片、队列指针 & 提供调度决策所需的全部依据 \\
内存管理 & 页表基址、段表、界限寄存器 & 描述地址空间并实施访问保护 \\
I/O 与资源 & 打开文件表、设备表、信号量 & 记录进程占用的资源与通信对象 \\
记账信息 & CPU 时间、缺页次数、启动时间 & 用于统计、计费与性能分析 \\
\bottomrule
\end{tabulary}
\end{table}

\subsection{切换时的保存与恢复}

\noindent 当进程离开 CPU 时，内核把「处理器现场」写入其 PCB；当下次被调度时，再从 PCB 恢复到处理器寄存器。核心是以下寄存器集合：

\begin{equation}
    \text{现场} = \{\, \text{PC},\ \text{通用寄存器},\ \text{PSW},\ \text{栈指针} \,\}
\end{equation}

\noindent 其中 PC 决定「执行到哪条指令」，通用寄存器保存中间结果，PSW 保存标志位与特权级，栈指针决定函数调用与返回的位置。保存与恢复必须\textbf{成对且完整}，否则进程恢复后会出现结果错乱甚至崩溃。此外，PCB 中的状态、队列指针等字段也会被同步更新，以保证调度程序读到的是最新状态。

\section{上下文切换}

\noindent 上下文切换 (context switch) 是内核保存当前进程现场、调入下一个进程现场的过程。它把「一个进程停下、另一个进程开始」在硬件层面真正实现。切换本身不产生任何用户价值，属于\textbf{纯开销}，因此其效率直接影响系统性能。

\subsection{切换步骤}

\begin{enumerate}
    \item 保存当前进程 $P_0$ 的 PC、通用寄存器、PSW、栈指针到 $P_0$ 的 PCB。
    \item 更新 $P_0$ 的状态（运行 $\to$ 就绪，或运行 $\to$ 阻塞），并插入相应队列。
    \item 调用调度程序，按策略从就绪队列中选择下一个待运行进程 $P_1$。
    \item 恢复 $P_1$ 的 PCB 中保存的现场到处理器寄存器。
    \item 切换地址空间（更新页表基址寄存器、刷新 TLB），返回用户态继续执行。
\end{enumerate}

\noindent 第 1、4 步是现场读写，第 3 步是策略决策，第 5 步是地址空间与高速缓存相关的隐性代价。

\subsection{开销来源}

\begin{itemize}
    \item \textbf{寄存器保存/恢复}：读写 PCB 与寄存器组，需要若干条内核指令。
    \item \textbf{模式切换}：用户态与内核态之间的特权级往返。
    \item \textbf{地址空间与缓存}：切换页表、失效 TLB，并污染 cache，使新进程初期命中率下降。
    \item \textbf{调度开销}：遍历就绪队列、计算优先级的策略成本。
\end{itemize}

\begin{equation}
    T_{\text{ctx}} = T_{\text{save}} + T_{\text{dispatch}} + T_{\text{restore}}
\end{equation}

\noindent 其中 $T_{\text{save}}$、$T_{\text{restore}}$ 为寄存器与现场读写开销，$T_{\text{dispatch}}$ 为调度选择开销。切换开销与硬件支持密切相关，通常在微秒量级。在时间片轮转中，若时间片为 $q$、一次切换为 $T_{\text{ctx}}$，则 CPU 的\textbf{有效利用率}与\textbf{切换开销占比}分别为

\begin{equation}
    U = \frac{q}{q + T_{\text{ctx}}},
    \qquad
    \text{开销占比} = \frac{T_{\text{ctx}}}{q + T_{\text{ctx}}} \times 100\%
\end{equation}

\noindent 时间片越短，切换越频繁，开销占比越高；因此必须在响应时间与切换开销之间折中。

\subsection{数值例题：上下文切换开销}

\noindent 设某系统采用时间片轮转调度，时间片 $q = 5\,\text{ms}$，一次上下文切换平均耗时 $T_{\text{ctx}} = 0.2\,\text{ms}$。试计算单个时间片的切换开销占比，并分析 $q$ 变化的影响。

\medskip
\noindent 由式 (6)：
\[
    \text{开销占比} = \frac{0.2}{5 + 0.2} \times 100\% = \frac{0.2}{5.2} \times 100\% \approx 3.85\%
\]
\[
    U = \frac{5}{5.2} \approx 96.15\%
\]

\noindent 即约 $3.85\%$ 的 CPU 时间被用于切换。若系统中同时有 $4$ 个就绪进程，则一个完整轮转周期为
$n(q + T_{\text{ctx}}) = 4 \times 5.2 = 20.8\,\text{ms}$，其中切换耗去 $4 \times 0.2 = 0.8\,\text{ms}$。
单位时间内发生的切换次数为
$N_{\text{sw}} = \dfrac{1000}{q + T_{\text{ctx}}} = \dfrac{1000}{5.2} \approx 192$ 次/秒。
下表对比不同时间片下的开销占比。

\begin{table}[H]
\centering
\caption{时间片 $q$ 对切换开销的影响（$T_{\text{ctx}} = 0.2\,\text{ms}$）}
\begin{tabulary}{\textwidth}{CCCC}
\toprule
时间片 $q$ (ms) & 开销占比 & 有效利用率 $U$ & 切换次数/秒 \\
\midrule
1  & $16.7\%$ & $83.3\%$ & 833 \\
5  & $3.85\%$ & $96.15\%$ & 192 \\
10 & $1.96\%$ & $98.04\%$ & 98 \\
50 & $0.40\%$ & $99.60\%$ & 20 \\
\bottomrule
\end{tabulary}
\end{table}

\noindent 结论：把 $q$ 从 $5\,\text{ms}$ 降到 $1\,\text{ms}$，有效利用率从 $96.15\%$ 跌到 $83.3\%$——时间片过短会让 CPU 忙于「换人」而非「干活」。反之 $q$ 过大则响应变差，退化为先来先服务。

\section{调度队列与调度器}

\noindent 操作系统用一组\textbf{队列}组织所有进程，用\textbf{调度器}在队列之间搬运进程。理解队列结构是理解调度算法的基础。

\subsection{就绪队列与等待队列}

\begin{itemize}
    \item \textbf{就绪队列 (ready queue)}：所有处于就绪态、等待 CPU 的进程，通常以链表或优先队列实现。
    \item \textbf{等待队列 (wait queue)}：因等待 I/O 或某事件而阻塞的进程，通常按设备或事件类型各维护一个。
\end{itemize}

\noindent 进程在队列之间的迁移可用下图表示：被调度选中者进入 CPU；等待 I/O 者进入等待队列；I/O 完成后回到就绪队列。

\begin{figure}[H]
\centering
\begin{tikzpicture}[font=\small, >=stealth, every node/.style={align=center}]
    \node[draw, fill=green!12, minimum width=1.0cm] (r1) at ({0},{0}) {P1};
    \node[draw, fill=green!12, minimum width=1.0cm, right=0.18cm of r1] (r2) {P2};
    \node[draw, fill=green!12, minimum width=1.0cm, right=0.18cm of r2] (r3) {P3};
    \node[right=0.18cm of r3, font=\small] (dots) {$\cdots$};
    \node[above=0.22cm of r1, anchor=west, font=\scriptsize, green!45!black] {就绪队列};

    \node[draw, rounded corners, fill=orange!18, minimum width=2.6cm, below=1.4cm of r2] (sched) {调度器 / 分派};
    \draw[->] (r2) -- (sched);

    \node[draw, rounded corners, fill=blue!12, left=2.6cm of sched] (cpu) {CPU 运行};
    \draw[->] (sched) -- (cpu);

    \node[draw, fill=red!12, minimum width=1.0cm] (w1) at ({-0.2},{-4.2}) {P4};
    \node[draw, fill=red!12, minimum width=1.0cm, right=0.18cm of w1] (w2) {P5};
    \node[above=0.22cm of w1, anchor=west, font=\scriptsize, red!55!black] {等待队列（如磁盘 I/O）};

    \draw[->] (cpu) -- node[right, font=\scriptsize]{等待 I/O} (w1);
    \draw[->, dashed] (w2.north) to[bend left=35] node[right, font=\scriptsize]{事件完成} (r1.south);
\end{tikzpicture}
\caption{就绪队列、CPU 与等待队列之间的进程迁移}
\end{figure}

\subsection{调度器的层次}

\noindent 按作用范围与频率，调度器可分为三级。

\begin{table}[H]
\centering
\caption{三级调度器}
\begin{tabulary}{\textwidth}{LLL}
\toprule
调度器 & 职责 & 频率 \\
\midrule
长程调度（作业调度） & 从后备作业中选择进入内存，决定多道程序度 & 最低 \\
中程调度（交换调度） & 决定进程换入/换出内存，缓解内存压力 & 中等 \\
短程调度（CPU 调度） & 从就绪队列中选择下一个占用 CPU 的进程 & 最高 \\
\bottomrule
\end{tabulary}
\end{table}

\noindent 短程调度器即通常所说的「CPU 调度器」，它每次上下文切换都可能被触发，直接决定交互响应与吞吐。多数系统并不实现完整的三级调度，但概念上区分它们有助于理解内存与 CPU 资源的耦合。

\section{线程}

\noindent 进程既是资源分配单位又是执行单位，导致「若要并发，就必须为每个执行流单独分配地址空间」，代价高昂。\textbf{线程 (thread)} 把执行流从资源分配中解耦：同一进程内的多个线程共享地址空间与资源，却各自独立执行，从而以更小的开销获得更强的并发能力。线程是 \textbf{CPU 调度的基本单位}，进程是 \textbf{资源分配的基本单位}。

\subsection{进程与线程的对比}

\begin{table}[H]
\centering
\caption{进程与线程对比}
\begin{tabulary}{\textwidth}{LLL}
\toprule
维度 & 进程 & 线程 \\
\midrule
定位 & 资源分配单位 & 调度执行单位 \\
地址空间 & 相互独立 & 同进程内共享 \\
切换开销 & 大（含地址空间与 TLB 切换） & 小（仅寄存器、栈与 PC） \\
通信方式 & IPC、管道、共享内存 & 直接读写共享变量，快但需同步 \\
隔离性 & 一个进程崩溃不影响其他进程 & 一个线程崩溃可能拖垮整个进程 \\
拥有资源 & 文件、内存、设备 & 基本不拥有，借用所属进程资源 \\
私有内容 & 整个地址空间 & 栈、寄存器、PC、线程局部存储 \\
\bottomrule
\end{tabulary}
\end{table}

\subsection{线程的地址空间}

\noindent 下图把「共享」与「私有」两部分画在一起：同一进程的多个线程共享代码段、数据段与打开文件表，但各自拥有私有的栈、寄存器和 PC。

\begin{figure}[H]
\centering
\begin{tikzpicture}[font=\small, >=stealth, every node/.style={align=center}]
    \draw[rounded corners, fill=blue!3, blue!35] ({0},{0}) rectangle ({3.4},{4.4});
    \node[font=\scriptsize, blue!60!black] at ({1.7},{4.15}) {进程地址空间（共享）};
    \node[draw, fill=green!12, minimum width=2.8cm] (code) at ({1.7},{3.5}) {代码段};
    \node[draw, fill=green!12, minimum width=2.8cm] (data) at ({1.7},{2.8}) {数据段 / 堆};
    \node[draw, fill=green!12, minimum width=2.8cm] (files) at ({1.7},{2.1}) {打开文件表 / 信号};

    \node[draw, rounded corners, fill=orange!15, minimum width=3.6cm, minimum height=0.8cm] (t1) at ({7.2},{3.6}) {线程 1：PC / 寄存器 / 私有栈};
    \node[draw, rounded corners, fill=orange!15, minimum width=3.6cm, minimum height=0.8cm] (t2) at ({7.2},{2.45}) {线程 2：PC / 寄存器 / 私有栈};
    \node[draw, rounded corners, fill=orange!15, minimum width=3.6cm, minimum height=0.8cm] (t3) at ({7.2},{1.3}) {线程 3：PC / 寄存器 / 私有栈};

    \draw[<->] ({3.4},{3.5}) -- ({5.4},{3.5}) node[midway, above, font=\scriptsize]{共享};
    \draw[<->] ({3.4},{1.75}) -- ({5.4},{1.75}) node[midway, below, font=\scriptsize]{每线程私有};
\end{tikzpicture}
\caption{线程的地址空间：代码与数据共享，栈、寄存器、PC 每线程私有}
\end{figure}

\subsection{用户级线程与内核级线程}

\noindent 线程由谁管理，决定了它的可见性与并发能力。

\begin{table}[H]
\centering
\caption{用户级线程 (ULT) 与内核级线程 (KLT)}
\begin{tabulary}{\textwidth}{LLL}
\toprule
维度 & 用户级线程 ULT & 内核级线程 KLT \\
\midrule
管理者 & 用户态线程库 & 操作系统内核 \\
内核是否可见 & 不可见 & 可见，作为独立调度对象 \\
切换开销 & 很小，无需陷入内核 & 较大，需陷入内核 \\
多核并行 & 不能（一个 KLT 承载全部 ULT） & 可以，多个 KLT 可并行运行 \\
阻塞影响 & 一个线程阻塞导致整个进程阻塞 & 单个线程阻塞不影响其他线程 \\
典型实现 & 早期 \texttt{pthread} 的 N:1 & Linux 内核线程、NPTL \\
\bottomrule
\end{tabulary}
\end{table}

\subsection{多线程模型}

\noindent 用户级线程与内核级线程之间的映射方式即\textbf{多线程模型}。

\begin{table}[H]
\centering
\caption{多线程映射模型}
\begin{tabulary}{\textwidth}{LLL}
\toprule
模型 & 映射关系 & 特点 \\
\midrule
多对一 (many-to-one) & 多 ULT $\to$ 1 KLT & 切换快，但一线程阻塞即全进程阻塞，无法并行 \\
一对一 (one-to-one) & 1 ULT $\to$ 1 KLT & 并行好、阻塞隔离，但线程创建与切换开销大 \\
多对多 (many-to-many) & 多 ULT $\to$ 多 KLT & 折中，允许受限并行，实现复杂 \\
\bottomrule
\end{tabulary}
\end{table}

\noindent 一对一模型（如 Linux NPTL）实现简单、并行性好，但线程数受内核限制；多对多模型试图在灵活性与开销之间取得平衡，但内核与线程库的协作复杂，是实践中较少采用的折中方案。

\section{进程间通信简介}

\noindent 进程之间相互独立，若要让它们交换数据，就需要\textbf{进程间通信 (IPC)}。主流方式分为两大类：\textbf{共享内存}与\textbf{消息传递}。前者在地址空间中划出一块共享区域，进程直接读写，速度快但必须由用户自行同步；后者通过内核在进程之间搬移消息，天然同步、易于跨机器，但每次通信都要陷入内核。二者对比如下表。

\begin{table}[H]
\centering
\caption{共享内存与消息传递}
\begin{tabulary}{\textwidth}{LLL}
\toprule
维度 & 共享内存 & 消息传递 \\
\midrule
通信机制 & 映射同一物理内存，直接读写 & 以内核为中介收发消息 \\
速度 & 快（建立后无需系统调用） & 较慢（每次通信需陷入内核） \\
同步 & 用户负责，需信号量/互斥锁 & 发送/接收本身提供同步语义 \\
隔离性 & 弱，共享区出错会相互影响 & 强，仅通过消息交互 \\
是否跨主机 & 否 & 是（如 socket） \\
典型接口 & \texttt{shmget}/\texttt{mmap} & \texttt{pipe}、\texttt{msgqueue}、\texttt{socket} \\
\bottomrule
\end{tabulary}
\end{table}

\noindent 选择原则：\textbf{追求性能}且进程同机时用共享内存，并配合同步机制；\textbf{追求简单与安全}、或需要跨主机通信时用消息传递。同步与死锁的细节将在后续章节展开。

\section{本章小结}

\begin{itemize}
    \item 进程是程序的一次执行，由代码段、数据段和 PCB 组成；程序静态、进程动态。
    \item 进程有新建、就绪、运行、阻塞、终止五种状态；阻塞态必须先回到就绪态才能运行。
    \item PCB 是进程存在的唯一标志，保存标识、现场、调度、内存与资源等全部信息。
    \item 上下文切换由「保存 $P_0$ $\to$ 调度 $\to$ 载入 $P_1$」构成，是纯开销；开销占比 $= T_{\text{ctx}}/(q+T_{\text{ctx}})$。
    \item 就绪队列与等待队列配合三级调度器，完成进程在 CPU、内存与 I/O 之间的迁移。
    \item 线程是调度基本单位，进程是资源分配单位；同一进程的线程共享地址空间，私有栈与寄存器。
    \item 线程模型有多对一、一对一、多对多；用户级线程切换快但不能并行，内核级线程可并行但开销大。
    \item 进程间通信分共享内存与消息传递：前者快而需同步，后者安全且可跨主机。
\end{itemize}

\section{常见误区}

\begin{itemize}
    \item \textbf{把程序当进程}：程序是静态文件，进程是动态执行体；同一程序可对应多个进程。
    \item \textbf{认为阻塞可以直接变运行}：阻塞只回到就绪，必须再经调度才可能运行。
    \item \textbf{混淆就绪与阻塞}：就绪是「只差 CPU」，阻塞是「在等事件」，前者可被调度，后者不能。
    \item \textbf{把并发当并行}：单核只能并发；并行需要多核，二者不是同义词。
    \item \textbf{认为线程切换与进程切换等价}：线程切换不换地址空间，开销小得多。
    \item \textbf{认为多线程一定更快}：线程间共享数据需要同步，竞争激烈时并行收益可能被锁开销抵消。
    \item \textbf{认为用户级线程能利用多核}：一个内核线程承载全部用户级线程时，无法真正并行。
    \item \textbf{低估上下文切换成本}：时间片过短会使切换成为主要开销，CPU 有效利用率显著下降。
\end{itemize}
